\documentclass[11pt]{article}
% Préambule commun aux cours de la formation C++.
% Inclus par chaque cours.tex via % Préambule commun aux cours de la formation C++.
% Inclus par chaque cours.tex via % Préambule commun aux cours de la formation C++.
% Inclus par chaque cours.tex via \input{preambule}.

\usepackage[utf8]{inputenc}
\usepackage[T1]{fontenc}
\usepackage[french]{babel}
\usepackage{lmodern}
\usepackage{microtype}

\usepackage{geometry}
\geometry{a4paper, margin=2.5cm}

\usepackage{xcolor}
\usepackage{amsmath, amssymb}
\usepackage{enumitem}
\usepackage{hyperref}
\hypersetup{
    colorlinks=true,
    linkcolor=blue!60!black,
    urlcolor=blue!60!black,
    citecolor=blue!60!black,
}

% --- Coloration syntaxique du C++ via listings -----------------------------
\usepackage{listings}

\definecolor{codebg}{rgb}{0.97,0.97,0.95}
\definecolor{codekw}{rgb}{0.10,0.10,0.60}
\definecolor{codecomment}{rgb}{0.20,0.50,0.20}
\definecolor{codestring}{rgb}{0.65,0.15,0.15}
\definecolor{coderule}{rgb}{0.80,0.80,0.80}

% Permet les caractères accentués (UTF-8) à l'intérieur des blocs de code.
\lstset{
    literate=%
        {é}{{\'e}}1 {è}{{\`e}}1 {ê}{{\^e}}1 {ë}{{\"e}}1
        {à}{{\`a}}1 {â}{{\^a}}1 {ä}{{\"a}}1
        {î}{{\^i}}1 {ï}{{\"i}}1
        {ô}{{\^o}}1 {ö}{{\"o}}1
        {ù}{{\`u}}1 {û}{{\^u}}1 {ü}{{\"u}}1
        {ç}{{\c c}}1 {É}{{\'E}}1 {È}{{\`E}}1 {À}{{\`A}}1
}

\lstdefinestyle{cpp}{
    language=C++,
    backgroundcolor=\color{codebg},
    basicstyle=\ttfamily\small,
    keywordstyle=\color{codekw}\bfseries,
    commentstyle=\color{codecomment}\itshape,
    stringstyle=\color{codestring},
    showstringspaces=false,
    numbers=left,
    numberstyle=\tiny\color{gray},
    numbersep=8pt,
    frame=single,
    rulecolor=\color{coderule},
    breaklines=true,
    tabsize=4,
    morekeywords={constexpr,noexcept,override,final,nullptr,auto,
        std,string,vector,size_t,uint64_t,int64_t,optional,namespace}
}
\lstset{style=cpp}

% Raccourci pour du code en ligne, en police machine à écrire.
\newcommand{\code}[1]{\texttt{#1}}

% Boîte « À retenir »
\usepackage{tcolorbox}
\newtcolorbox{aretenir}[1][]{
    colback=yellow!8, colframe=orange!70!black,
    title=#1, fonttitle=\bfseries, sharp corners, boxrule=0.5pt
}
\newtcolorbox{piege}[1][Piège]{
    colback=red!5, colframe=red!60!black,
    title=#1, fonttitle=\bfseries, sharp corners, boxrule=0.5pt
}
.

\usepackage[utf8]{inputenc}
\usepackage[T1]{fontenc}
\usepackage[french]{babel}
\usepackage{lmodern}
\usepackage{microtype}

\usepackage{geometry}
\geometry{a4paper, margin=2.5cm}

\usepackage{xcolor}
\usepackage{amsmath, amssymb}
\usepackage{enumitem}
\usepackage{hyperref}
\hypersetup{
    colorlinks=true,
    linkcolor=blue!60!black,
    urlcolor=blue!60!black,
    citecolor=blue!60!black,
}

% --- Coloration syntaxique du C++ via listings -----------------------------
\usepackage{listings}

\definecolor{codebg}{rgb}{0.97,0.97,0.95}
\definecolor{codekw}{rgb}{0.10,0.10,0.60}
\definecolor{codecomment}{rgb}{0.20,0.50,0.20}
\definecolor{codestring}{rgb}{0.65,0.15,0.15}
\definecolor{coderule}{rgb}{0.80,0.80,0.80}

% Permet les caractères accentués (UTF-8) à l'intérieur des blocs de code.
\lstset{
    literate=%
        {é}{{\'e}}1 {è}{{\`e}}1 {ê}{{\^e}}1 {ë}{{\"e}}1
        {à}{{\`a}}1 {â}{{\^a}}1 {ä}{{\"a}}1
        {î}{{\^i}}1 {ï}{{\"i}}1
        {ô}{{\^o}}1 {ö}{{\"o}}1
        {ù}{{\`u}}1 {û}{{\^u}}1 {ü}{{\"u}}1
        {ç}{{\c c}}1 {É}{{\'E}}1 {È}{{\`E}}1 {À}{{\`A}}1
}

\lstdefinestyle{cpp}{
    language=C++,
    backgroundcolor=\color{codebg},
    basicstyle=\ttfamily\small,
    keywordstyle=\color{codekw}\bfseries,
    commentstyle=\color{codecomment}\itshape,
    stringstyle=\color{codestring},
    showstringspaces=false,
    numbers=left,
    numberstyle=\tiny\color{gray},
    numbersep=8pt,
    frame=single,
    rulecolor=\color{coderule},
    breaklines=true,
    tabsize=4,
    morekeywords={constexpr,noexcept,override,final,nullptr,auto,
        std,string,vector,size_t,uint64_t,int64_t,optional,namespace}
}
\lstset{style=cpp}

% Raccourci pour du code en ligne, en police machine à écrire.
\newcommand{\code}[1]{\texttt{#1}}

% Boîte « À retenir »
\usepackage{tcolorbox}
\newtcolorbox{aretenir}[1][]{
    colback=yellow!8, colframe=orange!70!black,
    title=#1, fonttitle=\bfseries, sharp corners, boxrule=0.5pt
}
\newtcolorbox{piege}[1][Piège]{
    colback=red!5, colframe=red!60!black,
    title=#1, fonttitle=\bfseries, sharp corners, boxrule=0.5pt
}
.

\usepackage[utf8]{inputenc}
\usepackage[T1]{fontenc}
\usepackage[french]{babel}
\usepackage{lmodern}
\usepackage{microtype}

\usepackage{geometry}
\geometry{a4paper, margin=2.5cm}

\usepackage{xcolor}
\usepackage{amsmath, amssymb}
\usepackage{enumitem}
\usepackage{hyperref}
\hypersetup{
    colorlinks=true,
    linkcolor=blue!60!black,
    urlcolor=blue!60!black,
    citecolor=blue!60!black,
}

% --- Coloration syntaxique du C++ via listings -----------------------------
\usepackage{listings}

\definecolor{codebg}{rgb}{0.97,0.97,0.95}
\definecolor{codekw}{rgb}{0.10,0.10,0.60}
\definecolor{codecomment}{rgb}{0.20,0.50,0.20}
\definecolor{codestring}{rgb}{0.65,0.15,0.15}
\definecolor{coderule}{rgb}{0.80,0.80,0.80}

% Permet les caractères accentués (UTF-8) à l'intérieur des blocs de code.
\lstset{
    literate=%
        {é}{{\'e}}1 {è}{{\`e}}1 {ê}{{\^e}}1 {ë}{{\"e}}1
        {à}{{\`a}}1 {â}{{\^a}}1 {ä}{{\"a}}1
        {î}{{\^i}}1 {ï}{{\"i}}1
        {ô}{{\^o}}1 {ö}{{\"o}}1
        {ù}{{\`u}}1 {û}{{\^u}}1 {ü}{{\"u}}1
        {ç}{{\c c}}1 {É}{{\'E}}1 {È}{{\`E}}1 {À}{{\`A}}1
}

\lstdefinestyle{cpp}{
    language=C++,
    backgroundcolor=\color{codebg},
    basicstyle=\ttfamily\small,
    keywordstyle=\color{codekw}\bfseries,
    commentstyle=\color{codecomment}\itshape,
    stringstyle=\color{codestring},
    showstringspaces=false,
    numbers=left,
    numberstyle=\tiny\color{gray},
    numbersep=8pt,
    frame=single,
    rulecolor=\color{coderule},
    breaklines=true,
    tabsize=4,
    morekeywords={constexpr,noexcept,override,final,nullptr,auto,
        std,string,vector,size_t,uint64_t,int64_t,optional,namespace}
}
\lstset{style=cpp}

% Raccourci pour du code en ligne, en police machine à écrire.
\newcommand{\code}[1]{\texttt{#1}}

% Boîte « À retenir »
\usepackage{tcolorbox}
\newtcolorbox{aretenir}[1][]{
    colback=yellow!8, colframe=orange!70!black,
    title=#1, fonttitle=\bfseries, sharp corners, boxrule=0.5pt
}
\newtcolorbox{piege}[1][Piège]{
    colback=red!5, colframe=red!60!black,
    title=#1, fonttitle=\bfseries, sharp corners, boxrule=0.5pt
}


\title{\textbf{Chapitre 2 — \code{string} \& \code{vector} au quotidien,
       et organiser son code}\\[0.3em]
       \large Les deux conteneurs de base, et la modularité}
\author{Formation C++ moderne}
\date{}

\begin{document}
\maketitle

\begin{abstract}
\noindent Au chapitre 1 tu as croisé \code{std::string} et \code{std::vector}
en survol. Ici on les prend pour de bon : ce sont les deux conteneurs que tu
manipuleras dans 90\% de ton code. On apprend leurs méthodes utiles, comment
les parcourir proprement (boucle par intervalle \emph{et} itérateurs en boîte
noire), comment classer un caractère sans tomber dans le piège classique de
\code{<cctype>}, et on découvre deux premières fonctions de \code{<algorithm>}.
La seconde moitié du chapitre passe à l'organisation : séparer
\code{.hpp}/\code{.cpp}, comprendre l'erreur « multiple definition », ranger son
code dans des \emph{namespaces}, et un rappel de RAII. Chaque brique est
présentée \emph{avant} que tu en aies besoin, avec sa signature et un exemple
d'appel.
\end{abstract}

\tableofcontents
\bigskip\hrule\bigskip

% ===========================================================================
\section{\code{std::string} en pratique}

\code{std::string} gère sa mémoire, connaît sa taille, se concatène avec
\code{+} et se compare avec \code{==}. C'est l'outil par défaut pour le texte ;
oublie les tableaux de \code{char} hérités du C. Voici les méthodes que tu
utiliseras le plus, chacune avec un exemple d'appel concret.

\subsection*{Taille : \code{size()}}

\code{s.size()} renvoie le nombre de caractères, sous la forme d'un
\code{std::size\_t} (un entier \emph{non signé}). \code{length()} est un
synonyme exact.

\begin{lstlisting}
std::string s = "salut";
std::size_t n = s.size();        // n vaut 5
bool vide = s.empty();           // false (equivaut a s.size() == 0)
\end{lstlisting}

\subsection*{Accès à un caractère : \code{at(i)} vs \code{[i]}}

Les deux donnent le \code{char} à l'indice \code{i} (à partir de 0). La
différence est la vérification des bornes :

\begin{lstlisting}
std::string s = "salut";
char a = s[0];       // 's'  -- PAS de verification de bornes
char b = s.at(4);    // 't'  -- verifie : leve std::out_of_range si i >= size()
char& r = s[0];      // reference : s[0] = 'S'; modifie la chaine en place
\end{lstlisting}

\begin{aretenir}[\code{at} ou \code{[]} ?]
\code{[]} est plus rapide (aucune vérification) mais un indice hors bornes est
un \emph{comportement indéfini} (UB) : crash, valeur corrompue, ou pire, rien
d'apparent. \code{at()} vérifie et lève une exception propre. En phase
d'apprentissage et partout où l'indice vient d'un calcul incertain, préfère
\code{at()}. Une fois sûr de tes bornes, \code{[]} est l'idiome courant.
\end{aretenir}

\subsection*{Concaténation : \code{+} et \code{+=}}

\begin{lstlisting}
std::string prenom = "Ada";
std::string nom    = "Lovelace";
std::string complet = prenom + " " + nom;   // "Ada Lovelace"
complet += " !";                            // "Ada Lovelace !"
complet += '?';                             // on peut aussi ajouter un seul char
\end{lstlisting}

\subsection*{Comparaison : \code{==}, \code{!=}, \code{<}}

Contrairement au C (où on devait appeler \code{strcmp}), les chaînes se
comparent directement. \code{<} suit l'ordre lexicographique.

\begin{lstlisting}
std::string a = "abc", b = "abc", c = "abd";
bool egal     = (a == b);    // true
bool avant    = (a < c);     // true  ("abc" vient avant "abd")
\end{lstlisting}

\subsection*{Sous-chaîne : \code{substr(pos, len)}}

\code{s.substr(pos, len)} renvoie une \emph{nouvelle} chaîne : \code{len}
caractères à partir de l'indice \code{pos}. Si \code{len} est omis (ou dépasse),
on prend jusqu'à la fin.

\begin{lstlisting}
std::string s = "anticonstitutionnel";
std::string m1 = s.substr(0, 4);    // "anti"
std::string m2 = s.substr(4, 3);    // "con"
std::string m3 = s.substr(4);       // "constitutionnel"  (jusqu'au bout)
\end{lstlisting}

\subsection*{Recherche : \code{find(c)} et \code{std::string::npos}}

\code{s.find(c)} renvoie l'indice de la \emph{première} occurrence de \code{c}
(un caractère ou une sous-chaîne). Si rien n'est trouvé, elle renvoie la valeur
spéciale \code{std::string::npos}.

\begin{lstlisting}
std::string s = "a=42;b=7";
std::size_t pos = s.find('=');           // 1
if (pos != std::string::npos) {
    std::string cle = s.substr(0, pos);  // "a"
}
std::size_t absent = s.find('z');        // == std::string::npos (non trouve)
\end{lstlisting}

\begin{piege}[\code{npos} n'est pas $-1$]
\code{std::string::npos} est la plus grande valeur d'un \code{std::size\_t}
(non signé). Ne teste \emph{jamais} \code{find(...) >= 0} (toujours vrai pour un
non-signé !) ni \code{find(...) == -1}. Le seul test correct est
\code{== std::string::npos} ou \code{!= std::string::npos}.
\end{piege}

\subsection*{Vers une API C : \code{c\_str()}}

Certaines fonctions anciennes attendent un \code{const char*} terminé par un
zéro. \code{c\_str()} te le fournit sans copie.

\begin{lstlisting}
std::string chemin = "data.txt";
std::FILE* f = std::fopen(chemin.c_str(), "r");   // l'API C veut un const char*
\end{lstlisting}

\begin{aretenir}[Le réflexe \code{std::string}]
Taille avec \code{size()}, accès sûr avec \code{at()}, assemblage avec
\code{+}/\code{+=}, comparaison avec \code{==}, découpe avec \code{substr},
recherche avec \code{find} (+ test \code{npos}). Tu couvres déjà l'immense
majorité des besoins texte.
\end{aretenir}

% ===========================================================================
\section{\code{std::vector} en pratique}

Le \code{std::vector} est le tableau dynamique standard : une suite d'éléments
contiguë en mémoire, qui grandit toute seule. C'est ton conteneur par défaut
dès que tu veux « une liste de choses ».

\subsection*{Construire et remplir : \code{push\_back}, \code{size}}

\begin{lstlisting}
std::vector<int> v;            // vide
v.push_back(10);               // {10}
v.push_back(20);               // {10, 20}
std::size_t n = v.size();      // 2
bool vide = v.empty();         // false

std::vector<int> w = {1, 2, 3};        // initialisation directe
std::vector<double> z(5, 0.0);         // 5 elements valant 0.0
\end{lstlisting}

\code{push\_back(x)} ajoute \code{x} à la fin. \code{size()} renvoie le nombre
d'éléments (un \code{std::size\_t}). Le \code{vector} gère lui-même sa mémoire :
quand il est plein, il en réalloue automatiquement (on verra le détail
— capacité, \code{reserve} — dans un chapitre ultérieur).

\subsection*{Accès : \code{[]} vs \code{at}}

Même distinction que pour \code{std::string} : \code{[]} ne vérifie pas les
bornes (UB si hors limites), \code{at()} vérifie et lève
\code{std::out\_of\_range}.

\begin{lstlisting}
std::vector<int> v = {10, 20, 30};
int a = v[1];        // 20  -- pas de verification
int b = v.at(2);     // 30  -- verifie les bornes
v[0] = 99;           // ecriture en place : v = {99, 20, 30}
\end{lstlisting}

\subsection*{Un vector de vectors : la grille}

Pour une matrice ou une grille 2D, l'outil naturel est un
\code{std::vector<std::vector<int>>} : un vecteur dont chaque élément est
lui-même un vecteur (une ligne).

\begin{lstlisting}
std::vector<std::vector<int>> grille = {
    {1, 2, 3},
    {4, 5, 6}
};
int x = grille[1][2];        // 6  : ligne 1, colonne 2
std::size_t lignes   = grille.size();      // 2
std::size_t colonnes = grille[0].size();   // 3
\end{lstlisting}

\paragraph{Construire une grille dimensionnée d'un coup.} Le constructeur
\code{std::vector(taille, valeur)} crée \code{taille} copies de \code{valeur}.
En l'imbriquant, on fabrique une grille \code{rows}$\times$\code{cols}
pré-remplie :

\begin{lstlisting}
int rows = 3, cols = 4, fill = 0;
std::vector<std::vector<int>> grid(rows, std::vector<int>(cols, fill));
// 3 lignes, chacune etant une copie d'un vector de 4 zeros
grid[2][3] = 7;              // tous les indices existent deja
\end{lstlisting}

\begin{aretenir}[Lis la construction de l'intérieur vers l'extérieur]
\code{std::vector<int>(cols, fill)} = une ligne de \code{cols} éléments valant
\code{fill}. Puis \code{std::vector<...>(rows, cette\_ligne)} = \code{rows}
copies de cette ligne. Résultat : une grille toute prête, sans aucun
\code{push\_back}.
\end{aretenir}

% ===========================================================================
\section{Parcourir un conteneur}

\subsection*{La boucle \code{for} par intervalle}

Pour visiter tous les éléments d'un \code{std::vector}, d'une
\code{std::string}\ldots, l'idiome moderne est la \textbf{boucle \code{for} par
intervalle} (\emph{range-based for}) :

\begin{lstlisting}
std::vector<int> v = {10, 20, 30};
for (int x : v) {                 // x prend 10, puis 20, puis 30
    std::cout << x << ' ';
}
\end{lstlisting}

Comme pour les arguments de fonction, la question clé est : \code{x} est-il une
\emph{copie} ou une \emph{référence} de l'élément ? Cela dépend de comment tu le
déclares :

\begin{lstlisting}
for (int x : v)          { x *= 2; }   // COPIE : v reste {10,20,30}
for (auto& x : v)        { x *= 2; }   // REFERENCE : v devient {20,40,60}
for (const auto& x : v)  { lire(x);  } // reference en lecture seule
\end{lstlisting}

\begin{aretenir}[Quel qualificatif choisir ?]
\begin{itemize}[noitemsep]
    \item \code{const auto\&} → \textbf{le défaut} : lire sans copier. Idéal
          pour les gros objets (\code{std::string}, sous-vecteurs\ldots).
    \item \code{auto\&} → quand tu veux \textbf{modifier} chaque élément en
          place.
    \item \code{auto} / par valeur → seulement pour les types minuscules
          (\code{int}, \code{char}\ldots) ou si tu veux explicitement une copie.
\end{itemize}
\end{aretenir}

\subsection*{La forme itérateur explicite (boîte noire)}

Derrière la boucle par intervalle se cachent des \textbf{itérateurs} : des
objets qui « pointent » sur un élément et savent avancer au suivant. Tu
croiseras leur forme explicite dans du code existant :

\begin{lstlisting}
for (auto it = v.begin(); it != v.end(); ++it) {
    std::cout << *it << ' ';   // *it = l'element pointe
}
\end{lstlisting}

\code{begin()} pointe sur le premier élément, \code{end()} \emph{juste après} le
dernier (sentinelle de fin), \code{++it} avance d'un cran, \code{*it} lit (ou
écrit) l'élément pointé.

\begin{aretenir}[Itérateurs : pour l'instant, boîte noire]
Tu n'as pas besoin de maîtriser les itérateurs aujourd'hui : la boucle par
intervalle suffit à presque tout. On les \emph{approfondira au chapitre
Pointeurs / Algorithmes}. Retiens juste la mécanique \code{begin}/\code{end},
\code{++it}, \code{*it} pour lire le code des autres.
\end{aretenir}

\subsection*{Parcourir à l'envers : \code{rbegin()} / \code{rend()}}

Pour parcourir de la fin vers le début, il existe des itérateurs \emph{inverses}
\code{rbegin()} (dernier élément) et \code{rend()} (juste avant le premier).
\code{++} les fait reculer dans la séquence d'origine.

\begin{lstlisting}
std::string s = "abc";
for (auto it = s.rbegin(); it != s.rend(); ++it) {
    std::cout << *it;          // affiche : c b a
}
\end{lstlisting}

\paragraph{L'astuce d'inversion en une ligne.} Beaucoup de conteneurs ont un
constructeur qui prend un \emph{intervalle} \code{[debut, fin)} sous forme de
deux itérateurs, et recopie tout ce qui est dedans. En lui passant la paire
\emph{inverse}, on obtient une copie à l'envers :

\begin{lstlisting}
std::string s = "Lovelace";
std::string inverse(s.rbegin(), s.rend());   // "ecalevoL"
\end{lstlisting}

\begin{aretenir}[Inverser une chaîne sans boucle]
\code{std::string(s.rbegin(), s.rend())} construit une nouvelle chaîne en lisant
\code{s} de la fin au début. C'est l'idiome le plus court pour inverser. Là
encore, le \emph{pourquoi} (itérateurs, constructeur d'intervalle) sera détaillé
plus tard ; sers-t'en comme d'une formule toute faite pour le moment.
\end{aretenir}

% ===========================================================================
\section{Classifier un caractère avec \code{<cctype>}}

L'en-tête \code{<cctype>} fournit des fonctions pour \emph{tester} et
\emph{transformer} un caractère. Les plus utiles :

\begin{center}
\begin{tabular}{ll}
\hline
\textbf{Fonction} & \textbf{Renvoie vrai si\ldots / fait\ldots} \\
\hline
\code{std::isalpha(c)} & \code{c} est une lettre \\
\code{std::isdigit(c)} & \code{c} est un chiffre \code{'0'}--\code{'9'} \\
\code{std::isspace(c)} & \code{c} est une espace, tab, retour ligne\ldots \\
\code{std::tolower(c)} & renvoie la version minuscule de \code{c} \\
\code{std::toupper(c)} & renvoie la version majuscule de \code{c} \\
\hline
\end{tabular}
\end{center}

Exemple d'appels (attention à la ligne suivante sur le cast, c'est le point
crucial) :

\begin{lstlisting}
char c = 'A';
bool lettre = std::isalpha(static_cast<unsigned char>(c));   // true
bool chiffre = std::isdigit(static_cast<unsigned char>(c));  // false
char minus = static_cast<char>(std::tolower(static_cast<unsigned char>(c)));
// minus vaut 'a'
\end{lstlisting}

\begin{piege}[Le cast \code{static\_cast<unsigned char>(c)} est OBLIGATOIRE]
Ces fonctions sont héritées du C et attendent un \code{int} dont la valeur tient
dans \code{unsigned char} (ou vaut \code{EOF}). Or sur beaucoup de plateformes
\code{char} est \emph{signé} : un caractère accentué ou tout octet $\geq 128$
devient \emph{négatif}. Passer directement un \code{char} négatif est un
\textbf{comportement indéfini}. La parade systématique : envelopper l'argument
dans \code{static\_cast<unsigned char>(c)}. À prendre comme un réflexe
automatique chaque fois que tu appelles une fonction de \code{<cctype>}.
\end{piege}

\paragraph{Normaliser une lettre.} Mettre un caractère en minuscule, par
exemple pour comparer sans tenir compte de la casse :

\begin{lstlisting}
// renvoie c en minuscule (c suppose etre une lettre)
char en_minuscule(char c) {
    return static_cast<char>(std::tolower(static_cast<unsigned char>(c)));
}
\end{lstlisting}

Note le double \code{static\_cast} : un en \emph{entrée} (pour respecter le
contrat de \code{tolower} et éviter l'UB), un en \emph{sortie} (car
\code{tolower} renvoie un \code{int} qu'on veut reranger dans un \code{char}).

% ===========================================================================
\section{Découvrir \code{<algorithm>} en douceur}

L'en-tête \code{<algorithm>} contient une grande boîte à outils de fonctions qui
travaillent sur des \emph{intervalles} (un \code{begin} et un \code{end}). On en
fera tout un chapitre plus tard ; pour l'instant, juste deux fonctions très
parlantes, à utiliser en boîte noire.

\subsection*{Compter : \code{std::count}}

\textbf{Signature (simplifiée) :}
\begin{lstlisting}
std::count(debut, fin, valeur);   // -> combien d'elements egaux a `valeur`
\end{lstlisting}

\code{std::count(v.begin(), v.end(), valeur)} compte combien d'éléments de
l'intervalle \code{[begin, end)} sont égaux à \code{valeur}.

\begin{lstlisting}
std::vector<int> v = {3, 1, 3, 7, 3};
long n = std::count(v.begin(), v.end(), 3);   // 3 occurrences

std::string s = "mississippi";
long s_count = std::count(s.begin(), s.end(), 's');   // 4
\end{lstlisting}

\subsection*{Chercher : \code{std::find}}

\textbf{Signature (simplifiée) :}
\begin{lstlisting}
std::find(debut, fin, valeur);    // -> iterateur sur la 1re occurrence, ou `fin`
\end{lstlisting}

\code{std::find} renvoie un \emph{itérateur} sur la première occurrence ; s'il
ne trouve rien, il renvoie l'itérateur de fin (\code{end}). D'où le test
habituel contre \code{end()} :

\begin{lstlisting}
std::vector<int> v = {10, 20, 30};
auto it = std::find(v.begin(), v.end(), 20);
if (it != v.end()) {
    std::cout << "trouve a la position " << (it - v.begin()) << '\n';  // 1
}
\end{lstlisting}

\begin{aretenir}[Boîte noire assumée]
Tu peux te servir de \code{std::count} et \code{std::find} dès maintenant sans
comprendre leur implémentation : ce sont « compter / chercher sur un
intervalle ». Le \emph{pourquoi} des itérateurs et la richesse de
\code{<algorithm>} (transformations, tris, prédicats\ldots) seront le sujet du
\emph{chapitre Algorithmes}.
\end{aretenir}

% ===========================================================================
\section{Organiser : headers (\code{.hpp}) vs sources (\code{.cpp})}

Dès qu'un projet dépasse un fichier, on sépare \textbf{déclaration} et
\textbf{définition}. On \emph{déclare} dans le header (le « contrat » : quels
noms existent), on \emph{définit} dans le source (le corps).

\begin{lstlisting}
// geometry.hpp
#pragma once                 // evite la double inclusion
namespace geom {
    double rectangle_area(double w, double h);   // declaration (pas de corps)
}
\end{lstlisting}
\begin{lstlisting}
// geometry.cpp
#include "geometry.hpp"
namespace geom {
    double rectangle_area(double w, double h) { return w * h; }  // definition
}
\end{lstlisting}

\paragraph{\code{\#pragma once}.} Place cette ligne en tête de chaque header.
Elle garantit que le contenu n'est inclus qu'\emph{une fois} par unité de
traduction, même si plusieurs fichiers s'incluent en cascade. (L'alternative
historique est le trio \code{\#ifndef} / \code{\#define} / \code{\#endif} ;
\code{\#pragma once} est plus court et largement supporté.)

\begin{piege}[Erreur d'édition de liens : « multiple definition »]
Si tu \emph{définis} une fonction (non \code{inline}) directement dans un header
inclus par plusieurs \code{.cpp}, chaque \code{.cpp} en produit une copie, et le
linker voit plusieurs définitions du même symbole : il refuse avec
\emph{multiple definition of `geom::rectangle\_area'}. La règle simple :
\textbf{déclarations dans le \code{.hpp}, définitions dans le \code{.cpp}}.
(Exceptions : fonctions \code{inline}, templates, \code{constexpr} — qu'on verra
plus tard.)
\end{piege}

\begin{aretenir}[Le partage du travail]
Le \code{.hpp} dit \emph{ce qui existe} (à inclure par les utilisateurs) ; le
\code{.cpp} dit \emph{comment ça marche} (compilé une seule fois). Cette
séparation est aussi ce qui permet de recompiler vite : changer un corps dans un
\code{.cpp} ne force pas à tout reconstruire.
\end{aretenir}

% ===========================================================================
\section{Namespaces}

Les \textbf{namespaces} organisent le code et évitent les collisions de noms
(deux bibliothèques peuvent définir un \code{area} sans se gêner). Tout ce qui
vient de la bibliothèque standard vit dans \code{std::}.

\begin{lstlisting}
namespace geom {
    double rectangle_area(double w, double h);
}
// appel depuis l'exterieur : on qualifie avec le nom du namespace
double a = geom::rectangle_area(3.0, 4.0);
\end{lstlisting}

\paragraph{Rouvrir un namespace.} Un namespace n'est pas un bloc unique : on
peut le « rouvrir » dans plusieurs fichiers ou plusieurs endroits. Tout ce qui
est entre \code{namespace geom \{ ... \}} appartient à \code{geom}, où que ce
bloc se trouve.

\begin{lstlisting}
namespace geom { double rectangle_area(double w, double h); }
// ... plus loin, ou dans un autre fichier ...
namespace geom { double circle_area(double r); }   // meme namespace, complete
\end{lstlisting}

\begin{piege}[\code{using namespace std;}]
Ne le mets \emph{jamais} dans un header (tu l'imposerais à tous ceux qui
l'incluent), et évite-le même en \code{.cpp}. Il importe \emph{tout}
\code{std} et provoque des collisions sournoises (ton \code{count} contre
\code{std::count}, ton \code{distance} contre \code{std::distance}\ldots).
Écrire \code{std::cout} est un peu plus long mais explicite et sans danger. Si
un nom revient sans cesse, préfère un \code{using} \emph{ciblé}
(\code{using std::string;}) plutôt que tout le namespace.
\end{piege}

% ===========================================================================
\section{RAII en survol}

\textbf{RAII} (\emph{Resource Acquisition Is Initialization}) : une ressource
est acquise dans le constructeur d'un objet et \emph{libérée automatiquement}
dans son destructeur, à la fin de la portée. C'est \emph{le} pilier de C++ : pas
de \code{finally}, pas d'oubli de libération — la destruction est
\textbf{déterministe}, elle a lieu à l'accolade fermante.

\begin{lstlisting}
{
    std::vector<int> v(1000);   // memoire allouee a la construction
}   // <-- ici v est detruit, memoire liberee automatiquement
\end{lstlisting}

C'est exactement pour ça que \code{std::string} et \code{std::vector} sont si
confortables : ils acquièrent leur mémoire à la création et la rendent seuls à
leur destruction. Tu n'écris ni \code{new} ni \code{delete}, tu ne « fermes »
rien à la main.

\begin{aretenir}[Pourquoi ça nous concerne déjà]
Tout ce que tu fais avec \code{string}/\code{vector} repose sur RAII sans que tu
aies à y penser. Écrire \emph{tes propres} classes RAII (donc un destructeur)
demande de connaître les classes : on le mettra en pratique au \emph{chapitre
Classes}. Pour l'instant, retiens l'idée : la fin de portée libère tout.
\end{aretenir}

\bigskip\hrule\bigskip
\noindent\textbf{À toi de jouer.} Les exercices du dossier \code{exercices/}
mettent en pratique, en difficulté croissante :

\begin{itemize}[noitemsep]
    \item \textbf{parcours-vector} — range-for, \code{const auto\&} /
          \code{auto\&}, accès par indice.
    \item \textbf{palindrome} — \code{at}/\code{[]}, comparaison, et l'astuce
          \code{rbegin}/\code{rend}.
    \item \textbf{geometry-module} — séparer \code{.hpp}/\code{.cpp},
          \code{\#pragma once}, namespace \code{geom}.
    \item \textbf{histogramme} — \code{std::vector}, \code{<cctype>},
          \code{std::count}.
    \item \textbf{mini-grille} — \code{vector<vector<int>>}, construction
          \code{Grid(rows, vector<int>(cols, fill))}, double boucle.
    \item \textbf{string-toolkit} — \code{substr}, \code{find}, \code{npos},
          normalisation de casse.
    \item \textbf{split-maison} — découper une chaîne sur un séparateur en un
          \code{vector<std::string>}.
\end{itemize}

Le \textbf{projet de synthèse} (\code{projet/}) assemble le tout dans un
générateur de \textbf{rapport texte} : lecture d'un texte, comptage de mots et
de caractères, mise en forme tabulée. Tu y combineras conteneurs, parcours,
\code{<cctype>}, \code{<algorithm>} et une organisation propre en modules.
Bon courage !

\end{document}
